% 5-anexa-b.tex starts here
% Keep \raggedbottom for this appendix too (same rationale as 5-anexa-a.tex):
% the floats here (tab:pan-stiva, tab:pan-tooling, fig:mean-vs-pan,
% tab:pan-vs-altele) use [htbp] so a float too tall for the current page
% drifts to the next and the following text flows up to fill the slack,
% instead of the [H] standoff that left page-bottom holes. \flushbottom
% would stretch any residual slack into subsection gaps rather than letting
% it settle at the page bottom, so \raggedbottom stays.
\raggedbottom

\section{Manifest Arhitectural \acrshort{pan}}\label{anexa:manifest}

Această anexă consolidează arhitectura \gls{pan} propusă în lucrare, sintetizând definiția, justificările tehnologice, comparația cu arhitecturi alternative și direcțiile de evoluție.
Spre deosebire de \secref{cap:preliminarii} (care narează drumul către \gls{pan}) și \secref{cap:implementare} (care detaliază implementarea concretă în \textit{\gls{kookio}}), prezenta anexă oferă o privire de ansamblu organizată tematic --- utilă fie pentru o lectură de referință după ce lucrarea a fost parcursă integral, fie ca punct de plecare pentru cititorul care vrea să înțeleagă \gls{pan} fără să citească întregul studiu de caz.

\subsection{Definiție și principii structurale}\label{subsec:pan-definitie}

\gls{pan} (\acrlong{pan}) este o arhitectură full-stack monolitică modulară pentru aplicații \gls{web} construite integral în ecosistemul \gls{typescript}.
Cele trei componente eponime --- \gls{prisma} ca \gls{orm} declarativ, \gls{analog} ca \gls{metaframework} \gls{angular} cu suport \acrshort{ssr} și \gls{nx} ca orchestrator de \gls{monorepo} --- sunt elementele constitutive în jurul cărora se organizează o stivă coerentă cu cinci principii fundamentale:

\begin{enumerate}
  \item \textbf{Limbaj unic de la capăt la capăt.} Întreaga aplicație --- \acrshort{ui}, server, validare, persistență --- este scrisă în \gls{typescript}. Această decizie elimină categoria de erori provenite din desincronizarea contractelor între limbaje diferite (de exemplu, Java pe server și \gls{typescript} pe client) și permite reutilizarea uniformă a tipurilor.

  \item \textbf{O singură sursă de adevăr pentru contractele de date.} Schema bazei de date declarată în \texttt{schema.prisma} este sursa primară din care se derivă, automat, atât clientul \gls{prisma} \gls{typescript}, cât și schemele de validare \gls{zod} (via \texttt{prisma-zod-generator}). Tipurile aplicaționale sunt apoi compuse din aceste scheme prin operatorii \texttt{.pick()}, \texttt{.omit()}, \texttt{.extend()} și \texttt{z.infer<>}. O modificare în \texttt{schema.prisma} se propagă, la compilare, prin toate straturile.

  \item \textbf{Topologie modulară aciclică, aplicată la \acrshort{ci}.} Codul este organizat într-un \gls{monorepo} \gls{nx} cu patru grupuri de biblioteci (\texttt{feature}, \texttt{prisma}, \texttt{server}, \texttt{shared}) și o relație de dependență strict aciclică. Regula \texttt{enforce-module-boundaries} respinge la integrarea continuă orice import care încalcă topologia, garantând că structura arhitecturală nu se degradează pe parcursul evoluției proiectului.

  \item \textbf{Absența unui strat de \textit{repository} intermediar.} Procedurile \gls{trpc} apelează direct clientul \gls{prisma} (\texttt{ctx.prisma}), fără un strat intermediar de tip \textit{repository} orientat-obiect. Această decizie reduce numărul de straturi de adaptare și forțează ca logica de afaceri să se exprime în termenii primitivelor expuse de \gls{prisma} --- un compromis explicit între simplitate și abstracție.

  \item \textbf{Randare pe server ca mod implicit.} Combinația \gls{angular} \(+\) \gls{analog} livrează \acrshort{ssr} drept comportament implicit pentru toate rutele, iar \acrshort{ssg} este folosit punctual pentru paginile cu conținut stabil. Astfel, prima impresie a utilizatorului depinde de un HTML populat încă de la prima cerere, important atât pentru indexarea \acrshort{seo}, cât și pentru perceperea performanței.
\end{enumerate}

\subsection{Stiva tehnologică detaliată}\label{subsec:pan-stiva}

Fiecare nivel al stivei \gls{pan} este materializat de o tehnologie concretă, aleasă în raport cu alternative consacrate.
Cele paisprezece niveluri sunt grupate în două tabele după responsabilitate: \tabref{tab:pan-stiva} acoperă componentele aplicative (frontend, server, persistență), iar \tabref{tab:pan-tooling} acoperă tooling-ul de dezvoltare și operare.

\begin{table}[htbp]
\centering
\caption{Componentele aplicative ale arhitecturii \gls{pan} și justificarea fiecărei alegeri tehnologice.}
\label{tab:pan-stiva}
\small
\begin{tabularx}{\textwidth}{llY}
\toprule
\textbf{Nivel} & \textbf{Tehnologie} & \textbf{Justificare} \\
\midrule
Frontend         & \gls{angular} 21                  & Tipizare strictă nativă, ecosistem Material/CDK consistent, semnale reactive \\
Meta-framework   & \gls{analog}                      & Singura opțiune nativ-\gls{angular} cu \acrshort{ssr}, rutare pe fișiere și \textit{endpoint}-uri server integrate \\
Server runtime   & \gls{nitro}                       & Adoptat de \gls{analog}, agnostic față de hosting (\gls{vercel}, Node.js, Bun, Deno) \\
API   & \gls{trpc}                        & Tipuri partajate de la capăt la capăt fără un strat suplimentar de definiții \\
Validare         & \gls{zod} via \texttt{prisma-zod-generator} & Generare automată din schema \gls{prisma}, integrare nativă cu \gls{trpc} \\
\gls{orm}        & \gls{prisma}                      & Schemă declarativă, client \gls{typescript} generat, migrații versionate \\
Bază de date     & \gls{cockroachdb}                 & Compatibilitate de protocol cu \gls{postgresql}, scalare orizontală nativă, nivel \textit{serverless} \\
Cache            & \gls{redis}                       & Suport nativ pentru \textit{rate-limiting}, structuri de date specializate \\
\bottomrule
\end{tabularx}
\end{table}

\begin{table}[htbp]
\centering
\caption{Tooling-ul de dezvoltare și operare al arhitecturii \gls{pan} și justificarea fiecărei alegeri tehnologice.}
\label{tab:pan-tooling}
\small
\begin{tabularx}{\textwidth}{llY}
\toprule
\textbf{Nivel} & \textbf{Tehnologie} & \textbf{Justificare} \\
\midrule
Monorepo         & \gls{nx}                          & Filiație \gls{angular} (Nrwl), cache distribuit, generatoare native, \texttt{enforce-module-boundaries} \\
Build            & \gls{vite}                        & Suport \gls{typescript} nativ, \textit{hot module replacement} rapid, compatibilitate directă cu \gls{vitest} \\
Test unitar      & \gls{vitest}                      & Rulează pe \gls{vite} (nu duplică configurația), API compatibil Jest \\
Test \acrshort{e2e} & \gls{playwright}               & Suport multi-browser nativ, paralelism prin \textit{workers}, depanare cu \acrshort{ui} dedicată \\
\acrshort{cicd}  & GitHub Actions cu \glspl{compositeaction} & Aliniere cu țintele \gls{nx}, reutilizare prin acțiuni compozite \\
Hosting          & \gls{vercel}                      & Suport nativ pentru \gls{analog}/\gls{nitro}, mediu de previzualizare per ramură, variabile de mediu centralizate \\
\bottomrule
\end{tabularx}
\end{table}

\subsection{Comparație cu arhitecturi alternative}\label{subsec:pan-comparatie}

\gls{pan} se înscrie într-o linie de arhitecturi full-stack moștenite conceptual din \gls{mean}, propusă de Karpov în 2013~\cite{karpovMean2013}.
\figref{fig:mean-vs-pan} ilustrează paralela structurală între cele două.

\begin{figure}[htbp]
  \centering
  % The PNG carries ~5% baked-in white padding on each side (content occupies
  % x[100..1898] of 2000px ≈ 89.9% width). Scaling to 1.11x makes the diagram's
  % *content* span the text block; the padding overflows symmetrically into the
  % 2.5cm page margin (~0.9cm each side, well within margin). \centering keeps
  % the overflow balanced. Adjust the factor if the source PNG is re-exported
  % with different padding.
  \makebox[\textwidth]{\includegraphics[width=1.11\textwidth]{mean-vs-pan}}
  \caption{Paralelă structurală între arhitectura \gls{mean} (Karpov, 2013) și arhitectura \gls{pan} propusă în prezenta lucrare. Diagrama evidențiază maturizarea ecosistemului: \gls{analog} absoarbe rolurile distribuite în \gls{mean} între \textit{Express} și \gls{angular}, iar \gls{nx} adresează preocuparea organizării \gls{monorepo} care, în era \gls{mean}, nu era încă subiect de discuție arhitecturală.}
  \label{fig:mean-vs-pan}
\end{figure}

Spre deosebire de \gls{mean}, care era un set de tehnologii relativ disjuncte unite prin convenție comunitară, \gls{pan} reflectă o filiație tehnologică directă: comunitatea Nrwl (\gls{nx}) este fondată de foști ingineri Google din echipa \gls{angular}, iar comunitatea \gls{analog} menține \gls{nx} ca platforma preferată de orchestrare.
Această convergență structurală reduce frecușul de integrare pe care \gls{mean} îl avea (fiecare dezvoltator \gls{mean} își configura propriul flux de \textit{build}).

Stiva contemporană React (MERN + \gls{nextjs} + \textit{Turborepo}) ocupă același spațiu de soluții ca \gls{pan}, dar într-un alt ecosistem.
\tabref{tab:pan-vs-altele} oferă o comparație directă.

\begin{table}[htbp]
\centering
\caption{Comparație între \gls{mean} (2013), stiva React modernă și arhitectura \gls{pan}.}
\label{tab:pan-vs-altele}
\small
\begin{tabularx}{\textwidth}{lYYY}
\toprule
\textbf{Aspect} & \textbf{\gls{mean} (2013)} & \textbf{MERN + \gls{nextjs} + Turborepo} & \textbf{\gls{pan}} \\
\midrule
Frontend         & \gls{angular} (vechi)   & React                  & \gls{angular} 21 \\
Meta-framework   & --- (Express)           & \gls{nextjs}           & \gls{analog} \\
Server runtime   & Node.js + Express       & Node.js + \gls{nextjs} & \gls{nitro} \\
API   & REST manual             & API routes + \gls{trpc} opțional & \gls{trpc} obligatoriu \\
Tipuri partajate & Nu                      & Parțial (App Router)   & End-to-end automat \\
Bază de date     & MongoDB                 & MongoDB sau \gls{postgresql} & \gls{cockroachdb} \\
\gls{orm}        & Mongoose                & \gls{prisma} sau Drizzle & \gls{prisma} \\
Monorepo         & ---                     & Turborepo              & \gls{nx} \\
Limbaj           & JavaScript              & \gls{typescript} parțial & \gls{typescript} pur \\
\bottomrule
\end{tabularx}
\end{table}

Diferențele majore între \gls{pan} și stiva React modernă sunt:

\begin{itemize}
  \item \textbf{Comunitate.} \gls{pan} este construit pe \gls{angular}, framework cu o cotă de utilizare mai mică decât a React (\(50{,}1\,\%\) față de \(81{,}1\,\%\) \textit{State of JavaScript 2024}~\cite{stateOfJS2024}). Această diferență este atât o limitare practică (mai puțini dezvoltatori disponibili, mai puține biblioteci de \acrshort{ui} \textit{third-party}), cât și o invitație la disciplină --- \gls{angular} impune o organizare mai structurată implicit, prin injectarea dependențelor și componente \textit{standalone} strict tipizate.

  \item \textbf{Maturitate.} \gls{nextjs} are aproape un deceniu de evoluție și o adopție extinsă în producție; \gls{analog} are sub trei ani și o comunitate mică. \gls{pan} compensează prin alegerea \gls{angular} și \gls{nitro} (ambele mature individual) ca platforme de bază, dar acceptă riscul comunitar al \gls{analog}.

  \item \textbf{Tipizare end-to-end.} \gls{trpc} este în \gls{pan} o componentă obligatorie a arhitecturii, nu opțională. Această decizie face \gls{pan} mai restrictiv decât stiva React modernă (unde poți alege REST, GraphQL sau \gls{trpc} ad-hoc), dar elimină categoria de erori de desincronizare contractuală între client și server.
\end{itemize}

Dincolo de comparația cu stivele full-stack contemporane, \gls{pan} se poziționează și față de bifurcația arhitecturală mai largă dintre monolit modular și microservicii. Un câștig recunoscut pe scară largă al microserviciilor este delimitarea contextelor (\textit{bounded contexts}): Soldani et al.~\cite{soldani2018} îl identifică drept câștigul de arhitectură primar (24 de studii), sursa guvernanței descentralizate, a toleranței la defecte și a flexibilității de adăugare și eliminare a componentelor. Arhitectura \gls{pan} reproduce acest beneficiu de modularitate prin feliile verticale \texttt{libs/feature/*} și prin regula \texttt{enforce-module-boundaries}, dar îl mută din planul rulării în planul compilării: topologia aciclică este verificată la integrarea continuă, iar \texttt{libs/feature/X} nu poate importa din \texttt{libs/feature/Y} (\secref{subsec:topologie-biblioteci}). Astfel, \gls{pan} capturează câștigul de design al contextelor delimitate fără a plăti durerile operaționale ale distribuției --- localizarea serviciilor, coordonarea și eșecurile în cascadă, pe care același studiu le enumeră ca dureri de management. În același registru, o parte dintre pattern-urile pe care literatura le propune ca mitigări ale durerilor de microservicii --- \textit{service discovery}, \textit{circuit breaker}, \textit{API gateway} --- sunt non-probleme într-un monolit modular; singura excepție notabilă este \textit{database per service}, față de care \gls{pan} adoptă conștient soluția inversă (o bază partajată), schimbând izolarea datelor pe simplitatea tranzacțiilor ACID native.

\subsection{Direcții de evoluție arhitecturală}\label{subsec:pan-evolutie}

\gls{pan}, ca arhitectură, este extensibilă pe mai multe axe care nu au fost explorate în studiul de caz \textit{\gls{kookio}}:

\begin{itemize}
  \item \textbf{Comunicare în timp real.} \gls{trpc} suportă \textit{subscriptions} prin WebSocket sau Server-Sent Events; o ramură naturală de evoluție este integrarea cu un broker de mesaje (\gls{redis} Pub/Sub, NATS) pentru notificări către clienți autentificați.

  \item \textbf{Funcții la nivelul rețelei de margine (\textit{edge functions}).} \gls{nitro} permite publicare parțială la nivelul \textit{edge} (\gls{vercel} Edge Functions, Cloudflare Workers). Procedurile \gls{trpc} care nu necesită \gls{prisma} pot fi mutate la \textit{edge} pentru reducerea latenței geografice.

  \item \textbf{Multi-regional.} \gls{cockroachdb} suportă nativ replicare cross-zone; pentru o aplicație care iese din regimul mono-regional, această capacitate poate fi activată fără modificări de cod aplicativ.

  \item \textbf{Aplicații mobile.} Stiva \gls{angular} are suport pentru aplicații mobile native prin Capacitor sau NativeScript; bibliotecile \texttt{shared} și \texttt{prisma} ale unui proiect \gls{pan} sunt direct reutilizabile între o aplicație \gls{web} și una mobilă.

  \item \textbf{Internaționalizare.} \gls{angular} i18n combinat cu rutare multilingvă \gls{analog} (rute parametrizate prin prefix de limbă) acoperă cazul aplicațiilor multilingve fără modificări arhitecturale.
\end{itemize}

\gls{pan} nu este, în concluzie, o arhitectură universal aplicabilă.
Este o arhitectură pentru echipe mici cu cunoștințe \gls{angular} consolidate care construiesc aplicații \gls{typescript} moderne, monolitice (în sensul publicării unitare) și centrate pe date.
În acest spațiu de soluții, \gls{pan} oferă o coerență tehnologică și o disciplină structurală pe care alternativele bazate pe React le ating doar prin alegeri și convenții \textit{ad-hoc}.
% 5-anexa-b.tex ends here
